\documentclass[11pt,twocolumn]{article}

\usepackage[margin=0.85in]{geometry}
\usepackage[T1]{fontenc}
\usepackage{amsmath,amssymb}
\usepackage{booktabs}
\usepackage{array}
\usepackage{hyperref}
\usepackage{xcolor}
\usepackage{microtype}
\usepackage{authblk}
\usepackage[numbers]{natbib}
\usepackage{enumitem}
\setlist{nosep}

\hypersetup{colorlinks=true,linkcolor=blue!60!black,citecolor=blue!60!black,urlcolor=blue!60!black}

\title{Constraining the Standard Codon-to-Amino-Acid Assignment\\by WToken Fine Classes and Position-Weighted Cost}

\author[1]{Jonathan Washburn}
\affil[1]{Independent Researcher}

\date{March 2026}

\begin{document}
\maketitle

\begin{abstract}
We present a constrained derivation of the standard genetic code's
codon-to-amino-acid assignment. The 8-tick DFT basis of Recognition
Science generates exactly 20 semantic atoms (WTokens), which are placed
in bijection with the 20 canonical amino acids. The WToken
mode/phase/offset structure partitions the amino acids into 8 fine
candidate classes, giving an upper bound of $36{,}864$ fine-class-respecting
assignments. However, these fine classes do not refine the standard code's
degeneracy pattern. When both constraints are imposed simultaneously, the
admissible search space collapses to an exact 48 assignments, and the
standard code is the unique optimum on this 48-state space under the
position-weighted mutation cost of a companion paper. Both the 48-state
count and the strict optimality theorem are machine-checked in Lean~4.
This shows that the combination of WToken structure, observed degeneracy
pattern, and position-weighted mutation cost is sufficient to single out
the standard code from a tiny explicitly characterized residual space.
\end{abstract}

\section{Introduction}

A global-optimality result by itself does not explain why the standard
genetic code assigns \emph{these} 20 amino acids to \emph{these} codon
families. A derivation must do more than say that the code is good under a
cost function; it must also constrain the admissible assignment space.

Recognition Science provides exactly such a constraint. The 8-tick DFT
structure generates 20 neutral semantic atoms (WTokens), and the canonical
WToken$\,\leftrightarrow\,$amino-acid bijection partitions the amino acids
into eight fine classes. At first glance this seems to reduce the search
space dramatically: if assignments must respect the fine classes, then
$20!$ collapses to an upper bound of $2!^6 \times 4!^2 = 36{,}864$.

The crucial subtlety is that the fine classes do \emph{not} refine the
standard degeneracy pattern. For example, Val and Leu lie in the same fine
class but have 4-fold and 6-fold degeneracy respectively. Likewise Ser and
Thr, Lys and Arg, and the four-member ``special imaginary'' class cross
degeneracy boundaries. So the naive 36,864-state fine-class space is not
the right object for an exact codon-assignment derivation.

This paper isolates the correct object: the intersection of
\begin{enumerate}
\item the WToken fine-class constraint, and
\item the observed standard-code degeneracy pattern.
\end{enumerate}
That exact intersection has only 48 assignments, and we prove in Lean~4
that the standard code is the unique optimum on this 48-state space under
the position-weighted mutation cost from the companion global-optimality
paper~\cite{Paper2}.

\section{The WToken Constraint}

\subsection{The 20 WTokens}

The 8-tick DFT basis gives four mode families: modes 1+7, 2+6, 3+5, and
the Nyquist mode 4. Together with four phase levels and the real/imaginary
offset available only in mode 4, this yields exactly 20 WTokens:
\[
3 \times 4 + 1 \times 4 \times 2 = 12 + 8 = 20.
\]
In Lean, this is the theorem \texttt{wtoken\_count} in
\texttt{WTokenIso.lean}. The canonical amino acid type also has cardinality
20 (\texttt{amino\_acid\_count}), and the explicit map
\texttt{wtoken\_to\_amino} is proved bijective
(\texttt{wtoken\_amino\_bijection}).

\subsection{Fine Candidate Classes}

The canonical bijection induces eight fine classes:

\begin{center}
\small
\begin{tabular}{@{}lll@{}}
\toprule
Class & Members & WToken origin \\
\midrule
Aliphatic small & Gly, Ala & mode 1+7, $\phi$ 0--1 \\
Aliphatic branched & Val, Leu & mode 1+7, $\phi$ 2--3 \\
Polar hydroxyl & Ser, Thr & mode 2+6, $\phi$ 0--1 \\
Polar amide & Asn, Gln & mode 2+6, $\phi$ 2--3 \\
Charged negative & Asp, Glu & mode 3+5, $\phi$ 0--1 \\
Charged positive & Lys, Arg & mode 3+5, $\phi$ 2--3 \\
Aromatic (real) & His, Phe, Tyr, Trp & mode 4, $\tau = 0$ \\
Special (imaginary) & Pro, Cys, Met, Ile & mode 4, $\tau = 2$ \\
\bottomrule
\end{tabular}
\end{center}

This partition is formalized by \texttt{FineCandidateClass},
\texttt{aminoFineClass}, and \texttt{wtokenFineClass} in
\texttt{WTokenAssignmentBridge.lean}; the theorem
\texttt{canonical\_bridge\_respects\_fineClass} proves that the canonical
WToken$\,\leftrightarrow\,$amino map respects these classes.

The resulting fine-class-respecting search space has size at most
\[
2!^6 \times 4!^2 = 36{,}864,
\]
formalized by \texttt{fineClassSearchBound\_eq\_factorials} in
\texttt{NearUniqueness.lean}.

\subsection{Why the Naive Fine-Class Space Is Not Enough}

The eight fine classes are chemically plausible, but they do not line up
with the standard code's degeneracy classes. Lean explicitly proves the
basic counterexample:
\[
\texttt{aminoFineClass Val} = \texttt{aminoFineClass Leu}
\quad\text{but}\quad
\texttt{degeneracyClass Val} \neq \texttt{degeneracyClass Leu}.
\]
The theorem name is
\texttt{sameFineClass\_not\_sameDegeneracy} in
\texttt{WTokenFineClassOptimality.lean}.

This matters because any exact derivation of the standard code must respect
the observed degeneracy pattern. A fine-class-respecting assignment can swap
Val and Leu, but such a swap changes a 4-fold family into a 6-fold family
and is therefore not an admissible reassignment of the actual code.

\section{The Exact 48-State Intersection}

\subsection{Intersecting Fine Classes with Degeneracy}

Once we impose both constraints simultaneously, only four nontrivial blocks
remain:
\begin{itemize}
\item $\{\mathrm{Gly}, \mathrm{Ala}\}$ inside the 4-fold degeneracy class,
\item $\{\mathrm{Asn}, \mathrm{Gln}\}$ inside the 2-fold class,
\item $\{\mathrm{Asp}, \mathrm{Glu}\}$ inside the 2-fold class,
\item $\{\mathrm{His}, \mathrm{Phe}, \mathrm{Tyr}\}$ inside the 2-fold class.
\end{itemize}

All other amino acids are fixed because their fine classes cross degeneracy
boundaries. The exact intersection count is therefore
\[
2! \times 2! \times 2! \times 3! = 48.
\]

This is formalized in Lean by the finite type \texttt{FineDegChoice} and
the theorems
\begin{itemize}
\item \texttt{fineDegChoice\_count}
  (\texttt{Fintype.card FineDegChoice = 48}),
\item \texttt{fineDegChoice\_count\_factorials}
  (\texttt{2! × 2! × 2! × 3! = 48}),
\item \texttt{choiceAssignment\_respectsFineClass},
\item \texttt{choiceAssignment\_respectsDegeneracy}.
\end{itemize}

\subsection{Position-Weighted Cost on the 48-State Space}

We now evaluate the exact 48-state space under the same position-weighted
mutation cost used in the companion global-optimality paper~\cite{Paper2}.
That cost is defined in \texttt{PositionWeightedCost.lean} and uses the
integer weight vector
\[
(104,\,55,\,202,\,4131,\,4,\,25,\,1609,\,3870),
\]
corresponding to:
\begin{itemize}
\item hydrophobicity$^2$ at positions 1, 2, 3,
\item charge$^2$ at position 3,
\item side-chain volume$^2$ at positions 1, 2,
\item biosynthetic pathway distance at positions 1, 2.
\end{itemize}

The exact theorem proved in Lean is:

\begin{quote}
\textbf{Theorem}
(\texttt{standard\_strictly\_best\_on\_fineDegChoices}). For every
non-identity \texttt{FineDegChoice} $c$,
\[
\texttt{posStandardCost globalOptimalWeights}
<
\texttt{posTotalCost globalOptimalWeights (choiceAssignment c)}.
\]
\end{quote}

In words: among the 48 assignments that respect both the WToken fine
classes and the standard degeneracy pattern, the standard code is the
unique optimum.

\subsection{Interpretation}

This yields a clean two-stage derivation:
\begin{enumerate}
\item The 8-tick DFT structure forces 20 WTokens and an 8-class fine
  partition of the amino acids.
\item Intersecting that partition with the observed degeneracy pattern leaves
  exactly 48 admissible assignments.
\item The position-weighted mutation cost uniquely selects the standard code
  from those 48.
\end{enumerate}

The derivation is therefore conditional in an important sense: it explains
the standard code once the WToken structure and the observed degeneracy
pattern are both imposed. It does not derive the degeneracy pattern itself.

\section{Formal Verification}

The structural pieces of this paper are formalized in Lean~4
(version 4.27.0) with Mathlib:

\begin{itemize}
\item \texttt{WTokenIso.lean}: the 20 WTokens, the explicit inverse, and
  the theorem \texttt{wtoken\_amino\_bijection}.
\item \texttt{WTokenAssignmentBridge.lean}: the fine classes and the theorem
  \texttt{canonical\_bridge\_respects\_fineClass}.
\item \texttt{NearUniqueness.lean}: the upper bound
  \texttt{fineClassSearchBound = 36{,}864}.
\item \texttt{WTokenFineClassOptimality.lean}: the exact 48-state
  intersection, the non-refinement counterexample
  \texttt{sameFineClass\_not\_sameDegeneracy}, and the theorem
  \texttt{standard\_strictly\_best\_on\_fineDegChoices}.
\item \texttt{PositionWeightedCost.lean}: the position-weighted cost and the
  classwise certificates inherited from Paper~2.
\end{itemize}

The new module contains zero axioms and zero \texttt{sorry} assertions.
Unlike the earlier overbroad version of this paper, the present draft does
not rely on an invalid transfer from the 363,004 classwise search to the
naive 36,864 fine-class space. The exact 48-state theorem is internal to
Lean.

\section{Discussion}

\subsection{Relation to Coevolution and Stereochemical Hypotheses}

The derivation shares a structural feature with the coevolution
hypothesis~\cite{Wong1975}: the WToken fine classes group amino acids by
chemical and biosynthetic family, and the position-weighted cost gives
large weight to biosynthetic pathway distance. The derivation goes beyond
coevolution in two respects:
\begin{enumerate}
\item the classes are derived from the WToken DFT structure rather than
  being postulated from biochemistry, and
\item the exact 48-state residual space is fully enumerated and proved to
  have a unique optimum.
\end{enumerate}

It is less aligned with the stereochemical hypothesis~\cite{Woese1966},
since it does not posit direct codon--amino-acid affinity. The mechanism is
structural rather than chemical: the WToken organization constrains the
assignment space, and the mutation cost selects within that constrained
space.

\subsection{What Is and Is Not Derived}

The derivation explains:
\begin{itemize}
\item why there are exactly 20 amino acids (= 20 WTokens),
\item why those amino acids fall into the observed fine chemical families,
\item why the standard code is the unique optimum once both WToken fine
  classes and the observed degeneracy pattern are imposed.
\end{itemize}

It does not explain:
\begin{itemize}
\item why the degeneracy pattern itself is what it is,
\item why stop codons occupy their current positions,
\item the physical mechanism connecting WToken mode structure to amino acid
  chemistry and tRNA recognition.
\end{itemize}

\section{Conclusion}

The exact codon-to-amino-acid assignment problem is not solved by WToken
structure alone, but it is sharply constrained by it. The 8-tick DFT basis
forces 20 WTokens and an 8-class fine partition of the amino acids; when
that partition is intersected with the observed degeneracy pattern, only 48
assignments remain. The position-weighted mutation cost then uniquely
selects the standard code from this exact 48-state space.

This is a stronger and more precise result than the earlier overbroad claim
about a naive 36,864-state fine-class derivation. The corrected theorem is
fully formalized in Lean~4 and gives an exact constrained-derivation
statement for the standard code.

\section*{Data Availability}

All Lean source files are in the \texttt{reality} repository under
\texttt{IndisputableMonolith/Genetics/} and
\texttt{IndisputableMonolith/Water/WTokenIso.lean}.

\begin{thebibliography}{99}
\bibitem{FreelandHurst1998}
Freeland, S.J. and Hurst, L.D. (1998). The genetic code is one in a
million. \textit{J.~Mol.~Evol.} \textbf{47}(3), 238--248.

\bibitem{HaigHurst1991}
Haig, D. and Hurst, L.D. (1991). A quantitative measure of error
minimization in the genetic code. \textit{J.~Mol.~Evol.} \textbf{33}(5),
412--417.

\bibitem{Paper2}
Washburn, J. (2026). The standard genetic code is globally optimal under
position-weighted mutation cost. \textit{In preparation}.

\bibitem{Woese1966}
Woese, C.R., Dugre, D.H., Dugre, S.A., Kondo, M. and Saxinger, W.C.
(1966). On the fundamental nature and evolution of the genetic code.
\textit{Cold Spring Harb.~Symp.~Quant.~Biol.} \textbf{31}, 723--736.

\bibitem{Wong1975}
Wong, J.T.-F. (1975). A co-evolution theory of the genetic code.
\textit{Proc.~Natl.~Acad.~Sci.~USA} \textbf{72}(5), 1909--1912.
\end{thebibliography}

\end{document}
